\section{Cross-Task Analysis}\label{sec:cross-task}

Individually, each task's result is one comparison, with its own caveats. Together, they form
a pattern worth examining directly, without overstating any single piece of it.

\textbf{The similarity graph is close to a deterministic function of its own node features.}
The similarity edges used throughout this paper are constructed as $A_{\mathrm{sim}}=g(X)$
for a fixed threshold rule $g$ (Section~\ref{sec:data}). Because $A_{\mathrm{sim}}$ is a
deterministic function of $X$ and nothing else, it cannot, in principle, carry information
about a downstream target $Y$ beyond what $X$ already carries:
$I(Y;X,A_{\mathrm{sim}})=I(Y;X)$, a direct instance of the data-processing inequality
\cite{cover2006elements}. This is a claim about Shannon information -- that a sufficiently
powerful model, trained without limit, could not extract more from $(X,A_{\mathrm{sim}})$
jointly than from $X$ alone -- not that any particular finite, trained model here reaches
that limit; the experiments, not this inequality alone, test whether the trained models
benefit in practice. This covers the similarity edges specifically, not the deployed
construction's temporal edges, which encode segment order and are not recoverable from an
unordered feature set. It is not only theoretical: the deployed graphs, for the two datasets
this paper's comparisons use, have mean edge density 95.4\%-97.9\%, with 76.4\%-84.0\% of
individual graphs fully connected, and randomly rewiring the edge set reproduces almost the
same graph under a different name (Section~\ref{sec:data}).

\textbf{The direct graph-vs-edge-free comparisons are mixed at the single-cell level, but
none shows the graph clearly ahead once its own best supporting diagnostic is considered.}
Task 2's pooling-matched comparison favors the edge-free encoder by +0.0072 macro-F1 (a
larger, pooling-confounded +0.045 gap also exists, reported separately,
Section~\ref{sec:results-t2}). Task 3's standalone edge-free \textit{gnn\_only} comparison is
genuinely mixed and pooling-dependent -- the graph-equivalent baseline is slightly ahead at
mean-pool (0.1640 vs.\ 0.1599) and behind at max-pool (0.1640 vs.\ 0.1689), both validation
scores (Section~\ref{sec:results-t3}) -- so we do not read this as evidence in either
direction; Task 3's more informative, test-set evidence is branch-zeroing and
complementarity -- branch-zeroing finds a small positive contribution (+0.0088 macro-F1, not
zero, just small), while complementarity, corrected for the graph-only checkpoint's
label-polarity collapse (Section~\ref{sec:results-t3}), finds essentially none. Task 4's
edge-free and graph-based encoders show no reliable difference across three seeds on any
metric (Section~\ref{sec:results-t4}). Taken together: no task shows a clear graph advantage
across the available test-set evidence -- Task 2's control has disclosed dropout and seed
mismatches (Section~\ref{sec:protocol}), Task 3's edge-free comparison is validation-only,
and Task 4 supplies the repeated test-set comparison; where a standalone comparison is mixed
or pooling-dependent (Task 3), we say so rather than select the reading that fits the
pattern.

\textbf{The correlational diagnostic finds essentially no complementary signal; the causal
diagnostic finds a small but genuine one.} These are two different diagnostics, not one, and do not need to agree in degree to agree in
spirit. Complementarity (Section~\ref{sec:results-t3}) is correlational: comparing two
independently trained, single-branch models, the graph-only model is correct on 508 of
1{,}047 cells where text is wrong, a 48.5\% headline "rescue rate." This is misleading: the
graph-only checkpoint has exactly 0.0 precision and recall on 45 of 50 tags, so splitting the
508 by label polarity shows only 3 (0.6\%) are genuine true positives; the other 505 of 508
(99.4\%) are the graph model defaulting to "absent" while text wrongly said "present" -- an
always-predict-absent classifier would rescue essentially the same 48.6\%, so this number
does not distinguish the checkpoint from a trivial one. Branch-zeroing is a within-checkpoint
inference-time intervention, not a correlational count: inside the trained fusion model,
removing the graph branch costs 0.0088 macro-F1 and changes 1.66\% of individual
predictions -- small, but real (Section~\ref{sec:results-t3} notes what this diagnostic does
and does not establish about training-time causation). The two are not in tension:
complementarity, properly read, finds only three positive-label rescues -- too few for
meaningful complementary detection; branch-zeroing finds the fusion model relies on the graph
branch -- topology and node features together, since zeroing removes both, not
topology-isolating -- only slightly, not not at all. They are not two measurements of the
same quantity: complementarity speaks to whether the independently-trained graph-only model
detects tags text misses (largely not), while branch-zeroing speaks to the trained fusion
model's own reliance (slight, but real, not topology-specific). We read them together only as
pointing the same direction here, not converging on one number -- and
Section~\ref{sec:discussion} reports that a second implementation found substantially greater
fusion-model reliance and substantially more genuine rescue, an unresolved disagreement we do
not paper over here.

\textbf{Text dominance and graph weakness are separate findings, not one explanation for the
other.} Caption text leaks a substantial fraction of target labels directly (69.1\%
label-macro-averaged lexical exposure, measured for Task 1), explaining part of Task 1's text
strength independent of the graph; the sanitization control measuring this precisely
(Section~\ref{sec:results}) was only run for Task 1, so we do not claim the same figure or
result for Task 3 -- Task 3 uses the same MusicCaps captions, so the same leakage channel
plausibly exists there too, but this was not separately measured. The graph's weakness is
established separately, by the edge-free controls, not inferred from text's strength. Both
are true at once, reported as two distinct findings.

\textbf{Dataset integrity is a genuine open question for the small comparisons, though not
for the largest ones.} The one confirmed leakage effect (GTZAN's 7 cross-split exact-duplicate pairs, 3 of 150 test
tracks affected, Section~\ref{sec:data}) is small relative to Task 2's CNN-vs-GraphSAGE gap
(0.168) and its own larger, pooling-confounded edge-free comparison (0.045); three memorized
tracks are unlikely to move either by that much. It is not demonstrably too small, however,
to matter for Task 2's smaller, pooling-matched gap (0.0072): macro-F1 can move more than a
raw accuracy count suggests when errors concentrate within one genre, and a duplicated
validation-split track can also affect which checkpoint is selected. We do not correct the
split and rerun Task 2 here -- that would change every Task 2 number and is left as a direct
next step -- so we state this as an unresolved caveat on the smallest comparison
specifically, not the pattern as a whole. A companion audit's repaired factorial
\cite{tismir_audit} repairs the \emph{graph construction} (a separate density defect,
unrelated to this leakage) and finds the same qualitative pattern, but reuses the same GTZAN
split, so it does \emph{not} escape this leakage caveat -- not claimed as leakage-free
corroboration. Task 4 adds a different kind of independence: it uses MusicCaps (checked clean
of split-level ID leakage, Section~\ref{sec:data}) and a completely different task type, so at
least one leg of the pattern does not share GTZAN's leakage risk at all.

\textbf{The pattern holds across three structurally different learning settings, though not
via the identical measurement in each.} Single-label classification (Task 2), multi-label classification with a strong second
modality present (Task 3), and ranking under a contrastive objective with no class labels
(Task 4) are three different problems, evaluated with three different metrics; we do not
claim one uniform comparison recurs across all three. Task 2's direct comparison favors
edge-free. Task 3's own comparison is mixed (Section~\ref{sec:results-t3}); what recurs is
not "edge-free ahead" but the same underlying conclusion -- little detectable graph
benefit -- reached via an inference-time intervention and a correlational diagnostic instead.
Task 4 shows no reliable difference either way. The consistent thread is the
\emph{conclusion} (no reliable graph benefit detected), not an identically-shaped comparison,
more informative than any one task alone but read with the caveats above (validation-only
numbers and a mixed standalone result in Task 3, unresolved leakage for Task 2's smallest
comparison, two separate implementations across Tasks 1--3 versus Task 4) -- a consistent
signal, not four independent, airtight proofs of the same fact.
